\begin{figure*}[t]
\centering
\setcounter{subfigure}{0}
\begin{minipage}{\textwidth}
\centering
{\footnotesize
\textcolor{figAll}{\rule[2pt]{10pt}{1.2pt}} All (171)\quad
\textcolor{paperBlue}{\rule[2pt]{10pt}{0.7pt}} Easy (116)\quad
\textcolor{figGreen}{\rule[2pt]{4pt}{0.7pt}\hspace{2pt}\rule[2pt]{4pt}{0.7pt}} Moderate (28)\quad
\textcolor{paperOrange}{\rule[2pt]{6pt}{0.7pt}\hspace{2pt}\rule[2pt]{1pt}{0.7pt}} Difficult (27)\par}
\vspace{3pt}
\ESPanel{.2414\linewidth}{panels/fig10/a.pdf}{Mean quality}%
\hfill%
\ESPanel{.2414\linewidth}{panels/fig10/b.pdf}{Lower-tail quality}%
\hfill%
\ESPanel{.2414\linewidth}{panels/fig10/c.pdf}{Between-team variation}%
\hfill%
\ESPanel{.2414\linewidth}{panels/fig10/d.pdf}{Complete outputs}

\end{minipage}
\caption{Group size affects mean Q, lower-tail Q, between-team variation and output completion differently. Curves use the same 171 images: 116 easy, 28 moderate and 27 difficult. Q summaries use scorable outputs; complete-and-evaluable rates use all 43,776 attempts per size. P10 is the image-specific 10th percentile, averaged over images. Shading marks 3--6 and 20--30 people. The at-least-half threshold admits exactly half of an even-sized team: the required support stays at two votes from three to four people, then rises to three at five people, contributing to odd--even alternation.}
\label{fig:group-size-171}
\end{figure*}
